\subsection{Symplectic Classification and Transvections}

Throughout this subsection, symplectic geometries are finite-dimensional over a fixed arbitrary field $\mathbb{F}$.

\begin{theorem} \label{thm:symplectic-geometries-hyperbolic-classification} \hlt{Classification of Symplectic Geometries}\\
Let $\mathcal{V}$ be a symplectic geometry, $R=\operatorname{rad}(\mathcal{V})$, $n=\dim(\mathcal{V})$, and $r=\operatorname{rank}(\mathcal{V})$.
\begin{enumerate}[label=(\roman*)]
\setlength{\itemsep}{0pt}
\item If $R=\{\mathbf{0}\}$, then $\mathcal{V}$ is hyperbolic and $n$ is even.
\item $\mathcal{V}=R\odot\mathcal{H}$ for a hyperbolic subspace $\mathcal{H}\leq\mathcal{V}$, and $R$ is totally singular.
\item $r\in2\mathbb{Z}_{\geq0}$. If $r=2k$, then $\mathcal{V}$ has a basis $\beta$. For these choices, $M_{\beta}=X_{2k,n-2k}\defeq\operatorname{diag}(J_2^{\oplus k},0_{n-2k})$, where $J_2\defeq\left(\begin{smallmatrix}0&1\\-1&0\end{smallmatrix}\right)$.
\item For symplectic geometries $\mathcal{V},\mathcal{W}$, $\mathcal{V}\cong\mathcal{W}\Longleftrightarrow\dim\mathcal{V}=\dim\mathcal{W}\ \wedge\ \operatorname{rank}(\mathcal{V})=\operatorname{rank}(\mathcal{W})$.
\end{enumerate}
\begin{proof}
\begin{enumerate}[label=(\roman*)]
\setlength{\itemsep}{0pt}
\item[$(i)$] The claim holds for $n=0$. Suppose $n>0$ and fix $\mathbf{0}\neq\mathbf{x}\in\mathcal{V}$.\\
The vector $\mathbf{x}$ is isotropic by alternation and lies in a hyperbolic plane $\mathcal{H}_1$ by Theorem \ref{thm:isotropic-vector-sits-in-hyperbolic-plane}.\\
By Theorem \ref{thm:metric-orthogonal-complement-dimension}(ii),(iv), $\mathcal{V}=\mathcal{H}_1\odot\mathcal{H}_1^{\perp}$ and $\mathcal{H}_1^{\perp}$ is nonsingular.\\
The complement $\mathcal{H}_1^{\perp}$ is hyperbolic by induction on $n$, so $\mathcal{V}$ is hyperbolic and $n$ is even.
\item[$(ii)$] By Theorem \ref{thm:radical-splitting-metric-space}, $\mathcal{V}=R\odot\mathcal{S}$ for a nonsingular $\mathcal{S}$, and $R$ is totally singular.\\
The subspace $\mathcal{S}$ is hyperbolic by (i), so $\mathcal{H}\defeq\mathcal{S}$ has the required properties.
\item[$(iii)$] By (ii), $\mathcal{V}=R\odot\mathcal{H}$ with $\mathcal{H}$ hyperbolic. Since $\dim R=n-r$, $\dim\mathcal{H}=r$.\\
Write $\mathcal{H}$ as an orthogonal sum of $k$ hyperbolic planes; then $r=2k$.\\
A hyperbolic basis of $\mathcal{H}$ followed by a basis of $R$ is a basis $\beta$ of $\mathcal{V}$ satisfying $M_{\beta}=X_{2k,n-2k}$.
\item[$(iv)$] ($\Rightarrow$) Let $\sigma:\mathcal{V}\to\mathcal{W}$ be an isometry. Then $\dim\mathcal{V}=\dim\mathcal{W}$ by bijectivity.\\
If $\beta$ is a basis of $\mathcal{V}$, then $\sigma\beta$ is a basis of $\mathcal{W}$ and $M_{\sigma\beta}=M_{\beta}$ by form preservation; hence $\operatorname{rank}(\mathcal{V})=\operatorname{rank}(\mathcal{W})$.
\item[$(iv)$] ($\Leftarrow$) Suppose $\dim\mathcal{V}=\dim\mathcal{W}$ and $\operatorname{rank}(\mathcal{V})=\operatorname{rank}(\mathcal{W})$. By (iii), $\mathcal{V}$ and $\mathcal{W}$ have bases with the same matrix $X_{2k,n-2k}$.\\
The basis-to-basis isomorphism is an isometry by Proposition \ref{prop:metric-form-matrix-recovery-and-congruence}(i), so $\mathcal{V}\cong\mathcal{W}$.
\end{enumerate}
\hlt{Structural Commentary}: Radical splitting transfers classification to hyperbolic planes, whose block matrix determines isometry.\\
\hlt{Key Insight}: Every nonsingular alternating summand is exhausted by hyperbolic planes.
\end{proof}
\end{theorem}

\begin{definition} \label{def:symplectic-transvection} \hlt{Symplectic Transvection}\\
Let $\mathcal{V}$ be a nonsingular symplectic geometry, $\mathbf{0}\neq\mathbf{v}\in\mathcal{V}$, and $a\in\mathbb{F}$. Define $\tau_{\mathbf{v},a}:\mathcal{V}\to\mathcal{V}$ by $\tau_{\mathbf{v},a}(\mathbf{x})=\mathbf{x}+a\innerprod{\mathbf{x},\mathbf{v}}\mathbf{v}$.\\
This operator is the \hlt{symplectic transvection} with direction $\mathbf{v}$ and parameter $a$.
\end{definition}

\begin{proposition} \label{prop:symplectic-transvection-properties} \hlt{Basic Properties of Symplectic Transvections}\\
Let $\mathcal{V}$ be a nonsingular symplectic geometry and $\mathbf{0}\neq\mathbf{v}\in\mathcal{V}$. Fix $a,b\in\mathbb{F}$, $c\in\mathbb{F}^{\times}$, and $\sigma\in\operatorname{Sp}(\mathcal{V})$.
\begin{enumerate}[label=(\roman*)]
\setlength{\itemsep}{0pt}
\item $\tau_{\mathbf{v},a}\in\operatorname{Sp}(\mathcal{V})$;
\item $\tau_{\mathbf{v},a}=I_{\mathcal{V}}\Longleftrightarrow a=0$;
\item $\tau_{\mathbf{v},a}\tau_{\mathbf{v},b}=\tau_{\mathbf{v},a+b}$;
\item $\tau_{\mathbf{v},a}^{-1}=\tau_{\mathbf{v},-a}$;
\item $\sigma\tau_{\mathbf{v},a}\sigma^{-1}=\tau_{\sigma\mathbf{v},a}$;
\item $\tau_{c\mathbf{v},a}=\tau_{\mathbf{v},ac^2}$.
\end{enumerate}
\begin{proof}
\begin{enumerate}[label=(\roman*)]
\setlength{\itemsep}{0pt}
\item[$(i)$] Put $T\defeq\tau_{\mathbf{v},a}$, $\lambda(\mathbf{x})\defeq\innerprod{\mathbf{x},\mathbf{v}}$. Then $(\forall\,\mathbf{x},\mathbf{y}\in\mathcal{V})\,\innerprod{T\mathbf{x},T\mathbf{y}}-\innerprod{\mathbf{x},\mathbf{y}}=a\lambda(\mathbf{y})\lambda(\mathbf{x})-a\lambda(\mathbf{x})\lambda(\mathbf{y})=0$ by bilinearity and alternation.\\
$(\forall\,\mathbf{x}\in\mathcal{V})\,\tau_{\mathbf{v},-a}T(\mathbf{x})=\mathbf{x}=T\tau_{\mathbf{v},-a}(\mathbf{x})$ by alternation. Thus $T$ is an invertible form-preserving linear map, so $T\in\operatorname{Sp}(\mathcal{V})$.
\item[$(ii)$] ($\Rightarrow$) Suppose $\tau_{\mathbf{v},a}=I_{\mathcal{V}}$. If $a\neq0$, fix $\mathbf{x}$ satisfying $\innerprod{\mathbf{x},\mathbf{v}}\neq0$ by nonsingularity.\\
Then $\tau_{\mathbf{v},a}\mathbf{x}-\mathbf{x}=a\innerprod{\mathbf{x},\mathbf{v}}\mathbf{v}\neq\mathbf{0}$, a contradiction. Hence $a=0$.
\item[$(ii)$] ($\Leftarrow$) If $a=0$, then $\tau_{\mathbf{v},a}=I_{\mathcal{V}}$ by definition.
\item[$(iii)$] By alternation, $(\forall\,\mathbf{x}\in\mathcal{V})\,\tau_{\mathbf{v},a}\tau_{\mathbf{v},b}(\mathbf{x})=\mathbf{x}+(a+b)\innerprod{\mathbf{x},\mathbf{v}}\mathbf{v}=\tau_{\mathbf{v},a+b}(\mathbf{x})$. Hence $\tau_{\mathbf{v},a}\tau_{\mathbf{v},b}=\tau_{\mathbf{v},a+b}$.
\item[$(iv)$] $\tau_{\mathbf{v},a}^{-1}=\tau_{\mathbf{v},-a}$ by (ii),(iii) with $b=-a$.
\item[$(v)$] Fix $\mathbf{x}\in\mathcal{V}$. Then $\innerprod{\sigma^{-1}\mathbf{x},\mathbf{v}}=\innerprod{\mathbf{x},\sigma\mathbf{v}}$ by form preservation. Hence $\sigma\tau_{\mathbf{v},a}\sigma^{-1}(\mathbf{x})=\mathbf{x}+a\innerprod{\mathbf{x},\sigma\mathbf{v}}\sigma\mathbf{v}=\tau_{\sigma\mathbf{v},a}(\mathbf{x})$.
\item[$(vi)$] $(\forall\,\mathbf{x}\in\mathcal{V})\,\tau_{c\mathbf{v},a}(\mathbf{x})=\mathbf{x}+ac^2\innerprod{\mathbf{x},\mathbf{v}}\mathbf{v}=\tau_{\mathbf{v},ac^2}(\mathbf{x})$ by bilinearity.
\end{enumerate}
\hlt{Structural Commentary}: Alternation controls preservation and composition; nonsingularity detects the zero parameter.\\
\hlt{Key Insight}: The rank-one square vanishes, and the two mixed form terms cancel.
\end{proof}
\end{proposition}

\begin{theorem} \label{thm:symplectic-group-generated-by-transvections} \hlt{Symplectic Transvections Generate}\\
If $\mathcal{V}$ is nonsingular symplectic, then $\operatorname{Sp}(\mathcal{V})=\left\langle\tau_{\mathbf{v},a}:\mathbf{0}\neq\mathbf{v}\in\mathcal{V},\ a\in\mathbb{F}\right\rangle$.
\begin{proof}
\begin{enumerate}[label=(\roman*)]
\setlength{\itemsep}{0pt}
\item \hlt{Generated Subgroup}: Put $G\defeq\left\langle\tau_{\mathbf{v},a}:\mathbf{0}\neq\mathbf{v}\in\mathcal{V},\ a\in\mathbb{F}\right\rangle$. Then $G\leq\operatorname{Sp}(\mathcal{V})$ by Proposition \ref{prop:symplectic-transvection-properties}(i).
\item \hlt{Vector Transitivity}: Fix $\mathbf{0}\neq\mathbf{x},\mathbf{y}\in\mathcal{V}$. If $\innerprod{\mathbf{x},\mathbf{y}}\neq0$, then $\tau_{\mathbf{y}-\mathbf{x},\innerprod{\mathbf{x},\mathbf{y}}^{-1}}(\mathbf{x})=\mathbf{y}$.\\
Suppose $\innerprod{\mathbf{x},\mathbf{y}}=0$ and $\mathbf{y}=\lambda\mathbf{x}$ for $\lambda\in\mathbb{F}^{\times}$. Fix $\mathbf{z}$ satisfying $\innerprod{\mathbf{x},\mathbf{z}}\neq0$ by nonsingularity. Then $\innerprod{\mathbf{z},\mathbf{y}}=-\lambda\innerprod{\mathbf{x},\mathbf{z}}\neq0$.\\
If $\mathbf{x},\mathbf{y}$ are independent, choose $\mathcal{V}=\operatorname{span}(\mathbf{x},\mathbf{y})\oplus\mathcal{C}$ by Proposition \ref{prop:complements-exist-finite-dimensions}. Define $f\in\mathcal{V}'$ by $f(\alpha\mathbf{x}+\beta\mathbf{y}+\mathbf{c})=\alpha-\beta$ for $\mathbf{c}\in\mathcal{C}$.\\
Fix $\mathbf{z}$ satisfying $(\forall\,\mathbf{t}\in\mathcal{V})\,f(\mathbf{t})=\innerprod{\mathbf{t},\mathbf{z}}$ by Theorem \ref{thm:riesz-map-for-nonsingular-metric-spaces}. Then $\innerprod{\mathbf{x},\mathbf{z}}=\innerprod{\mathbf{z},\mathbf{y}}=1$ by the definition of $f$ and alternation.\\
In either case, put $c\defeq\innerprod{\mathbf{x},\mathbf{z}}$ and $d\defeq\innerprod{\mathbf{z},\mathbf{y}}$. Thus $c,d\neq0$, $\tau_{\mathbf{z}-\mathbf{x},c^{-1}}(\mathbf{x})=\mathbf{z}$, and $\tau_{\mathbf{y}-\mathbf{z},d^{-1}}(\mathbf{z})=\mathbf{y}$; $G$ is transitive on $\mathcal{V}\setminus\{\mathbf{0}\}$.
\item \hlt{Hyperbolic-Pair Transitivity}: Fix hyperbolic pairs $(\mathbf{u},\mathbf{v})$ and $(\mathbf{u}',\mathbf{v}')$.\\
Choose $\gamma\in G$ satisfying $\gamma\mathbf{u}=\mathbf{u}'$ by (ii), and put $\mathbf{w}\defeq\gamma\mathbf{v}$. Then $(\mathbf{u}',\mathbf{w})$ is hyperbolic. If $\mathbf{w}=\mathbf{v}'$, put $\delta\defeq I_{\mathcal{V}}$.\\
If $p\defeq\innerprod{\mathbf{w},\mathbf{v}'}\neq0$, put $\delta\defeq\tau_{\mathbf{v}'-\mathbf{w},p^{-1}}$. Then $\delta\mathbf{u}'=\mathbf{u}'$ and $\delta\mathbf{w}=\mathbf{v}'$ by direct evaluation. Suppose $\mathbf{w}\neq\mathbf{v}'$ and $\innerprod{\mathbf{w},\mathbf{v}'}=0$.\\
If $a\mathbf{u}'+b\mathbf{w}+c\mathbf{v}'=\mathbf{0}$, then $b+c=0$ and $a=0$ by pairing with $\mathbf{u}'$ and $\mathbf{w}$.\\
Hence $b(\mathbf{w}-\mathbf{v}')=\mathbf{0}$, so $a=b=c=0$ and $\mathbf{u}',\mathbf{w},\mathbf{v}'$ are independent. Choose $\mathcal{V}=\operatorname{span}(\mathbf{u}',\mathbf{w},\mathbf{v}')\oplus\mathcal{C}$ by Proposition \ref{prop:complements-exist-finite-dimensions}.\\
Define $f\in\mathcal{V}'$ by $f(a\mathbf{u}'+b\mathbf{w}+c\mathbf{v}'+\mathbf{d})=a+b-c$ for $\mathbf{d}\in\mathcal{C}$. Fix $\mathbf{z}$ satisfying $(\forall\,\mathbf{t}\in\mathcal{V})\,f(\mathbf{t})=\innerprod{\mathbf{t},\mathbf{z}}$ by Theorem \ref{thm:riesz-map-for-nonsingular-metric-spaces}.\\
Then $\innerprod{\mathbf{u}',\mathbf{z}}=\innerprod{\mathbf{w},\mathbf{z}}=\innerprod{\mathbf{z},\mathbf{v}'}=1$ by the definition of $f$ and alternation.\\
The transvections $\tau_{\mathbf{z}-\mathbf{w},1}$ and $\tau_{\mathbf{v}'-\mathbf{z},1}$ fix $\mathbf{u}'$. The identities $\tau_{\mathbf{z}-\mathbf{w},1}(\mathbf{w})=\mathbf{z}$ and $\tau_{\mathbf{v}'-\mathbf{z},1}(\mathbf{z})=\mathbf{v}'$ hold by direct evaluation.\\
Thus $\delta\defeq\tau_{\mathbf{v}'-\mathbf{z},1}\tau_{\mathbf{z}-\mathbf{w},1}$ fixes $\mathbf{u}'$ and maps $\mathbf{w}$ to $\mathbf{v}'$. In every case, $\delta\gamma\in G$ maps $(\mathbf{u},\mathbf{v})$ to $(\mathbf{u}',\mathbf{v}')$.
\item \hlt{Dimension Induction}: Fix $\sigma\in\operatorname{Sp}(\mathcal{V})$ and put $n\defeq\dim\mathcal{V}$. We induct on $n$. If $n=0$, then $\sigma=I_{\mathcal{V}}\in G$. Suppose $n>0$.\\
Choose a hyperbolic plane $\mathcal{H}=\operatorname{span}(\mathbf{u},\mathbf{v})$ by Theorem \ref{thm:symplectic-geometries-hyperbolic-classification}(i). By (iii), fix $\gamma\in G$ mapping $(\sigma\mathbf{u},\sigma\mathbf{v})$ to $(\mathbf{u},\mathbf{v})$; hence $\gamma\sigma|_{\mathcal{H}}=I_{\mathcal{H}}$.\\
By Theorem \ref{thm:metric-orthogonal-complement-dimension}(ii),(iv), $\mathcal{V}=\mathcal{H}\odot\mathcal{H}^{\perp}$ and $\mathcal{H}^{\perp}$ is nonsingular.\\
The subspace $\mathcal{H}^{\perp}$ is $\gamma\sigma$-invariant by form preservation and pointwise fixation of $\mathcal{H}$.\\
Thus $(\gamma\sigma)|_{\mathcal{H}^{\perp}}\in\operatorname{Sp}(\mathcal{H}^{\perp})$ and is a product of transvections by induction.\\
Each factor $\tau_{\mathbf{x},a}$ extends to the same transvection on $\mathcal{V}$. Since $\mathbf{x}\in\mathcal{H}^{\perp}$, it fixes $\mathcal{H}$ pointwise.\\
The ambient product agrees with $\gamma\sigma$ on both orthogonal summands. Therefore $\gamma\sigma\in G$ and $\sigma=\gamma^{-1}(\gamma\sigma)\in G$.
\end{enumerate}
\hlt{Structural Commentary}: Vector transitivity upgrades to pair transitivity; hyperbolic-plane induction then factors every isometry.\\
\hlt{Key Insight}: Pair transitivity isolates a fixed hyperbolic plane for dimension induction.
\end{proof}
\end{theorem}

\begin{proposition} \label{prop:symplectic-transformations-have-determinant-one} \hlt{Symplectic Transformations Have Determinant One}\\
If $\mathcal{V}$ is nonsingular symplectic, then $(\forall\,\sigma\in\operatorname{Sp}(\mathcal{V}))\,\det(\sigma)=1$.
\begin{proof}
Fix $\sigma\in\operatorname{Sp}(\mathcal{V})$. It is a product of transvections by Theorem \ref{thm:symplectic-group-generated-by-transvections}.\\
By determinant multiplicativity, it suffices to prove $\det(\tau_{\mathbf{v},a})=1$ for $\mathbf{0}\neq\mathbf{v}\in\mathcal{V}$ and $a\in\mathbb{F}$.\\
The transvection $\tau_{\mathbf{v},a}$ fixes $\mathbb{F}\mathbf{v}$ by alternation and induces $I_{\mathcal{V}/\mathbb{F}\mathbf{v}}$ by its defining formula.\\
Thus its matrix in a basis beginning with $\mathbf{v}$ is upper triangular with unit diagonal, so $\det(\tau_{\mathbf{v},a})=1$.\\
\hlt{Structural Commentary}: Transvection generation reduces the determinant to one basis-adapted rank-one map.\\
\hlt{Key Insight}: Only the leading coordinate changes, so the matrix is unit upper triangular.
\end{proof}
\end{proposition}
